% 灵息杯第一届正赛 —— 中文卷
% 请用 XeLaTeX 编译（需要 ctex 宏包，TeX Live / CTeX 套装均自带）
\documentclass[UTF8]{ctexart}
\usepackage{graphicx}
\usepackage{amsfonts}
\usepackage{amssymb}
\usepackage{amsmath}
\title{第一届灵息杯数学公开赛}
\author{夏一欢（XIA Yihuan）}
\date{\today}
\setlength{\parindent}{0pt}

\begin{document}

\maketitle

\section{引言}

欢迎参加第一届灵息杯数学公开赛！
比赛时间为 2026 年 10 月 1 日 00:00 至 2027 年 2 月 5 日（除夕）23:59（北京时间）。
全卷共 5 道题，每题 20 分，满分 100 分。

\paragraph{约定。}
全文 $\mathbb{N}=\{0,1,2,3,\ldots\}$——按中国教材的习惯，自然数\emph{包含}
$0$，即 $0\in\mathbb{N}$——
$\mathbb{Z}^{*}=\mathbb{Z}\setminus\{0\}$。
第 3 题用到记号 $K(S)$；没有学过代数学的选手请自行查阅``域扩张''这个词条。

\section{试题}
\subsection{第 1 题}

证明：对任意 $n,m\in\mathbb{N}\setminus\{0,1\}$（即 $n,m\ge 2$），存在
$x,y,z\in\mathbb{N}$，使得

$$x^n+y^n \equiv z^n \pmod{m}
\qquad\text{且}\qquad
z^n \not\equiv 0 \pmod{m}$$

（$x,y,z$ 须两两不同。注意 $0\in\mathbb{N}$，故它们都可以取 $0$。）

20 分。

\subsection{第 2 题}

设

$$f(x)=\sum_{i=0}^n a_i x^i\ \Bigl(\sum_{j=1}^n a_j^2\neq 0\Bigr)\in \mathbb{Z}[x]$$

证明：对任意常数 $C\in \mathbb{R}$，

$${\lvert\{u \in \mathbb{R}\mid f(u) = C\}\rvert\leq n}$$

20 分。

\subsection{第 3 题}

设

$$K_0=\mathbb{Q}$$

$$G_i=\{p\in \mathbb{C}\mid p^o=q^r,\ q\in K_i,\ r,o\in \mathbb{Z}^*\}$$

$$K_j=K_{j-1}(G_{j-1})$$

$$f(x)\in \mathbb{Z}[x]$$

并记 $n=\deg f$（全题中 $f$ 均假定为非常数）。``在 $K$ 上不可约''指 $f$
不能写成 $K[x]$ 中两个次数更小的多项式之积；因此只对非常数多项式谈论可约与否。

\subsubsection{(1)}

是否成立：($f(x)$ 在 $\mathbb{R}$ 上不可约) $\implies$ ($f(x)$ 在 $K_1$ 上
不可约)？回答并证明。

2 分。

\subsubsection{(2)}

是否成立：($f(x)$ 在 $K_1$ 上不可约) $\implies$ ($f(x)$ 在 $\mathbb{R}$ 上
不可约)？回答并证明。

6 分。

\subsubsection{(3)}

是否成立：$(\exists M \in \mathbb{N})(\forall f(x) \in \mathbb{Z}[x])[(f(x)$
在 $K_M$ 上不可约$) \implies (f(x)$ 在 $\mathbb{R}$ 上不可约$)]$？
（注意 $0\in\mathbb{N}$，故允许 $M=0$，此时 $K_M=K_0=\mathbb{Q}$。）

请就次数 $n=\deg f$ 指出该蕴含式对哪些 $n$ 成立、对哪些 $n$ 不成立，并证明。

6 分。

\subsubsection{(4)}

设

$$K^*_0=\mathbb{Q}$$

$$G^*_i=\{p\in \mathbb{R}\mid p^o=q^r,\ q\in K^*_i,\ r,o\in \mathbb{Z}^*\}$$

$$K^*_j=K^*_{j-1}(G^*_{j-1})$$

求一个 $f(x) \in \mathbb{Z}[x]$，使得它对一切 $M \in \mathbb{N}$ 都在
$K^*_M$ 上不可约，但它在 $\mathbb{R}$ 上可约，并给出它在 $\mathbb{R}$ 上的分解。

6 分。

\subsubsection*{注释}

(i) 次数 $\le4$ 的多项式方程都可用根式求解（求根公式、Cardano 公式、
Ferrari 公式）。

(ii)（Abel--Ruffini）次数 $\ge5$ 的多项式方程没有一般的根式解法；例如
$x^5-4x+2$ 的任一根都不能用根式表示。

(iii)（不可约情形 / casus irreducibilis）若 $\mathbb{Q}$ 上的不可约三次多项式
有三个实根，则它的任一根都不能用实根式表示。

\subsection{第 4 题}

国际象棋中马（knight）的八种走法可写成集合

$$K_2=\{\pm1\pm 2i,\ \pm 2 \pm i\}\subset\mathbb{Z}[i]$$

现定义``$a$-马''（strange A-knight）的走法集合为

$$K_a=\{\pm1\pm ai,\ \pm a \pm i\}\subset\mathbb{Z}[i]$$

（其中 $\pm1\pm ai$ 表示两个符号各自独立取值所得的四个数。）

\subsubsection{(1)}

证明：

$$a \equiv 1 \pmod{2}\implies\sum_{i=1}^n k_i\neq 1,
\quad\text{对任意 } n\ge 1 \text{ 与任意 } k_1,\ldots,k_n\in K_a$$

6 分。

\subsubsection{(2)}

证明：

$$a \equiv 0 \pmod{2}\implies
\bigl(\forall z\in\mathbb{Z}[i]\bigr)\ \bigl(\exists n\ge 1,\ 
k_1,\ldots,k_n\in K_a\bigr):\ \sum_{i=1}^n k_i=z$$

6 分。

\subsubsection{(3)}

设整数 $m\ge 4$，棋盘

$$B_m=\{a+bi\in\mathbb{Z}[i]\mid 0\le a,b<m\}$$

证明：对任意 $z\in B_m$，存在 $n\ge 1$ 与 $k_1,\ldots,k_n\in K_2$，使得

$$\sum_{i=1}^n k_i=z\qquad\text{且}\qquad
\sum_{i=1}^j k_i\in B_m\ \ \text{对一切 } 1\le j\le n ;$$

也就是说，普通的马从角上的 $0$ 出发，可以走遍棋盘的每一格，且全程不离开棋盘。

8 分。

\subsection{第 5 题}

设 $n\ge 1$，$K=(a_1,a_2,\ldots ,a_n)$ 是一个由 $0$ 和 $1$ 组成的有限序列。
按下列规则依次处理各个位置，得到新序列 $T(K)=(b_1,\ldots ,b_n)$。

\paragraph{规则。}
依次处理位置 $i=1,2,\ldots ,n$。轮到位置 $i$ 时，先按\emph{当前}的序列读出
$a_i$，按下面的条款决定是否改写；若需改写，则\emph{立刻}执行，然后才处理下一个
位置。因此，后面的计算会读到前面已经改过的值。以下各式中出现的 $a_j$ 一律指
处理时刻的当前值。
\begin{itemize}
\item 若当前 $a_i=1$，令 $p_i=\bigl|\{j<i\mid a_j=1\}\bigr|$ 为它前面 $1$ 的
个数。若 $p_i\ge 2$，则把第 $i+p_i-1$ 位改成 $1$。
\item 若当前 $a_i=0$，令 $s_i=\bigl|\{j>i\mid a_j=0\}\bigr|$ 为它后面 $0$ 的
个数。若 $s_i\ge 2$，则把第 $i-s_i+1$ 位改成 $0$。
\end{itemize}
其余情形什么都不做：数出的个数为 $0$ 或 $1$ 时不做；所指的位置不在
$\{1,\ldots ,n\}$ 内时也不做（不循环、不越界）。处理完位置 $n$ 后所得的序列
就是 $T(K)=(b_1,\ldots ,b_n)$。

\paragraph{问题。}
是否对任意 $n\ge 1$ 与任意 $K$，$T(K)$ 都必有\emph{不动点}，即存在
$t\in\{1,\ldots ,n\}$ 使得 $b_t=a_t$（$a_t$ 指\emph{原}序列 $K$ 的第 $t$ 位）？
若是，请给出证明；若否，请举出反例。

20 分。

\section{致谢与参赛须知}

赛题由夏一欢命制，由夏一欢与若干智能体试做。感谢 Loveinterpretation 不可或缺
的贡献。
允许组队，但每队只能提交一份答卷。在充分理解题意之后，请勿使用计算机、计算器、
大模型或书籍。（大模型与机器也可以提交答卷，但不参与评奖；若其解答优于参考
解答，将被附于参考解答之后。）

解题所需的一切知识都写在题后的注释中。我们相信各位都已至少学完初中数学。

请将解答以 PDF 形式发送至 aether\_math@163.com。手写与排版印刷的数学公式均可，
仅用纯文本而无规范记号书写的不予接收。请确保所有公式清晰可辨且数学上正确。

申诉与建议可发往同一邮箱。

\vspace{1em}
\noindent \textbf{祝好运！}

\end{document}
